% Stand 17. Juli 2026
%\subsection{Porträtfoto: Vorhersage durch lineare Regression}
%\label{AnhangPorträtfotoVorhersageRegression}
%\begin{figure}[H] 
%\centering
%\includegraphics[scale = 0.6]{Illustration_Kriging 40Prz Ergebnis OLS m Interaktion} 
%\caption{\small{Die \glqq Realität\grqq\ (links); durch Stichproben-Erhebungen gewonnene Kenntnis dieser
%\glqq Wahrheit\grqq\ (mittig); aus dem Halbwissen mittels lin.\,Regression
%gewonnene Vorhersage (rechts)}}
%\label{Abb_Illustration_Kriging 40Prz Ergebnis OLS m Interaktion}
%\end{figure} 
%
%Abb.\,(\ref{Abb_Illustration_Kriging 40Prz Ergebnis OLS m Interaktion}) zeigt die vergleichsweise deutlich weniger
%aussagekräftigen Ergebnisse einer multiplen linearen 
%Regression unter Berücksichtigung von Interaktionen, bei der 
%im Wesentlichen nur die globalen Mittelwerte einfließen und
%nachbarschaftliche Ähnlich\-keits\-be\-zie\-hun\-gen unberücksichtigt 
%bleiben. 
%
%\newpage

\subsection{Weitere Kovarianz-Funktionen}
\label{AnhangWeitereKovarianzFunktionen}
$\circ$ \underline{Stable Modell}
\begin{equation}
\rho_{stb}(r;\ell) :=  \exp\left( -\left( \frac{r}{\ell}\right)^\alpha \right),\ \alpha 
\in (0,\,2],
\end{equation}
wobei $\ell \in \mathbb{R}_+ \setminus \lbrace 0 \rbrace$.

\begin{figure}[H] 
\centering
\includegraphics[scale = 0.44]
{TheoModelle_Stabil_drei_verschiedene.png} 
\caption{\small{Korrelations-Funktion des 
stabilen Modells für verschiedene Parameter $\ell$}}
\label{Abb_StabilesModell_dreiLängen}
\end{figure} 

\begin{figure}[H] 
\centering
\includegraphics[scale = 0.5]{TheoModelle_Stabil_Zufallsfeld.png} 
\caption{\small{Realisierung eines isotropen, schwach stationären
Gauß'schen Zufallsfelds mit stabilem Modell (links) und das zugehörige Korrelogramm (rechts)}}
\label{Abb_StabilesModell_GRF}
\end{figure} 

Das Stable Modell spezifiziert wiederum eine ganze Modell Familie.
Insbesondere gilt (vgl.\,z.\,B.\,\cite[S.\,25]{Montero2015}):
\begin{itemize}[topsep=2pt,parsep=1pt,itemsep=0.5pt]
\item[$\bullet$] für $\alpha=1$ entspricht das Stable Modell dem 
exponentiellen Modell;
\item[$\bullet$] für $\alpha=2$ entspricht das Stable Modell dem 
Gauß'schen Modell.
\end{itemize}

%$\circ$ \underline{Exponentielles Modell}
%\begin{equation}
%\label{Gl_ExponentiellesModell}
%\rho_{exp}(r; \ell) :=  \exp\left( -\frac{r}{\ell} \right),\ 
%\ell \in \mathbb{R}_+ \setminus \lbrace 0 \rbrace
%\end{equation}
%
%\begin{figure}[H] 
%\centering
%\includegraphics[scale = 0.44]
%{TheoModelle_Exponentiell_drei_verschiedene.png} 
%\caption{\small{Korrelations-Funktion des 
%exponentiellen Modells für verschiedene Parameter $\ell$}}
%\label{Abb_ExponentiellesModell_dreiLängen}
%\end{figure} 
%
%\begin{figure}[H] 
%\centering
%\includegraphics[scale = 0.5]{TheoModelle_Exponentiell_Zufallsfeld.png} 
%\caption{\small{Realisierung eines isotropen, schwach stationären
%Gauß'schen Zufallsfelds mit exponentiellem Modell (links) und das zugehörige Korrelogramm (rechts)}}
%\label{Abb_ExponentiellesModell_GRF}
%\end{figure} 

\vspace*{10pt}

$\circ$ \underline{Gauß'sches Modell}
\begin{equation}
\label{Gl_GaussModell}
\rho_{gau}(r; \ell) := \exp\left( -\frac{r^2}{\ell^2} \right),\ \ell \in \mathbb{R}_+ \setminus \lbrace 0 \rbrace
\end{equation}

\begin{figure}[H] 
\centering
\includegraphics[scale = 0.44]
{TheoModelle_Gauss_drei_verschiedene.png} 
\caption{\small{Korrelations-Funktion des 
Gauß'schen Modells für verschiedene Parameter $\ell$}}
\label{Abb_GaussModell_dreiLängen}
\end{figure} 

\begin{figure}[H] 
\centering
\includegraphics[scale = 0.5]{TheoModelle_Gauss_Zufallsfeld.png} 
\caption{\small{Realisierung eines isotropen, schwach stationären
Gauß'schen Zufallsfelds mit Gauß'schen Modell (links) und das zugehörige Korrelogramm (rechts)}}
\label{Abb_GaussModell_GRF}
\end{figure} 


$\circ$ \underline{Matérn Modell}
\begin{equation}
\rho_{mat}(r; \ell, \nu ) := \frac{1}{\Gamma(\nu)2 ^{\nu -1}}
\left( \frac{\sqrt{2\nu}}{\ell} r \right)^\nu
\mathrm{K}_\nu \left( \frac{\sqrt{2\nu}}{\ell} r \right),
\end{equation}
wobei 
\begin{itemize}[topsep=2pt,parsep=1pt,itemsep=0.5pt]
\item[$\bullet$] $\Gamma(\cdot)$ die Gamma-Funktion,
\item[$\bullet$] $\mathrm{K}_\nu(\cdot)$ die hyperbolische Bessel-Funktion 
zweiter Art $\rightsquigarrow$ \href{https://en.wikipedia.org/wiki/Bessel_function#Modified_Bessel_functions}
{Wiki: Bessel function}
\item[$\bullet$] $\ell,\nu \in \mathbb{R}_+ \setminus \lbrace 0 \rbrace$
 Modell-Parameter
\end{itemize}
bezeichnen.\\

\begin{figure}[H] 
\centering
\includegraphics[scale = 0.44]
{TheoModelle_Matern_drei_verschiedene.png} 
\caption{\small{Korrelations-Funktion des 
Matérn-Modells für verschiedene Parameter $\ell$}}
\label{Abb_MaternModell_dreiLängen}
\end{figure} 

\begin{figure}[H] 
\centering
\includegraphics[scale = 0.5]{TheoModelle_Matern_Zufallsfeld.png} 
\caption{\small{Realisierung eines isotropen, schwach stationären
Gauß'schen Zufallsfelds mit Matérn Modell (links) und das zugehörige Korrelogramm (rechts)}}
\label{Abb_MaternModell_GRF}
\end{figure} 

Im Limes $\nu \rightarrow \infty$
konvergiert das Matérn Modell gegen das Gauß'sche
Modell, das heißt, 
das Gauß'sche Modell ist Teil der Matérn Modell Familie (vgl.\,z.\,B.\,\cite[S.\,84f]{Rasmussen2006}).\\

\newpage

\subsection{Ergänzende Analysen der Datenauswertung}
\label{AnhangErgänzendeAnalysenDatenauswertung}
\begin{figure}[H] 
\centering
\includegraphics[scale = 0.45]{Datenauswertung_P25MessungenApr2026DirektionaleSemivarianzBeispielFit340_NoConstraints} 
\caption{\small{Anpassung der Kovarianz-Funktion 
(\textcolor{NavyBlue}{\rule[1mm]{3mm}{0.5mm}}) an die
direktionale ($\vartheta=340^\circ$) empirische Semivarianz 
(\textcolor{Sepia}{$\bullet$});
oberhalb der empirischen Semivarianz sind die Umfänge 
des Lags vermerkt; der obere Wert der Beschränkung für den Range 
beträgt $100\,km$ }}
\label{Abb_Datenauswertung_P25MessungenApr2026DirektionaleSemivarianzBeispielFit340_NoConstraints}
\end{figure} 

\begin{figure}[H] 
\centering
\includegraphics[scale = 0.45]{Datenauswertung_P25MessungenApr2026VariogrammRoseRange_NoConstraints} 
\caption{\small{Direktionale Variogramm-Karte für den 
Range basierend auf den gemittelten 
$\mathrm{P}_{2.5}$- Messwerten für April 2026; der obere Wert der Beschränkung für den Range 
beträgt $100\,km$  }}
\label{Abb_Datenauswertung_P25MessungenApr2026VariogrammRoseRange_NoConstraints}
\end{figure} 

\begin{figure}[h]
\begin{center}
\includegraphics[width=0.44\textwidth]{Datenauswertung_P25_Apr2026_VariogrammKarte_Sill_Matern}
\includegraphics[width=0.44\textwidth]{Datenauswertung_P25_Jul2025_VariogrammKarte_Sill_Matern}
\end{center}
\caption{\small{Direktionale Variogramm-Karte für den 
Sill basierend auf den gemittelten 
$\mathrm{P}_{2.5}$- Messwerten für April 2026 (links) und Juli 2025 (rechts)}}
\label{Abb_Datenauswertung_P25_Apr2026Jul2025_VariogrammKarte_Sill_Matern}
\end{figure}


\begin{figure}[H] 
\centering
\includegraphics[scale = 0.45]{Datenauswertung_P25_LOO-CVNov2025_Ansicht_e_j} 
\caption{\small{Beispielhafte Darstellung der logarithmierten Fehler $e_j$ eines 
Durchlaufs der LOO-CV; zusatzlich sind der Mittelwert (\textcolor{WildStrawberry}{\rule[1mm]{3mm}{0.5mm}}) 
und die Standardabweichung (\textcolor{BurntOrange}{\rule[1mm]{3mm}{0.5mm}})
eingezeichnet}}
\label{Abb_Datenauswertung_P25_LOO-CVNov2025_Ansicht_e_j}
\end{figure} 


\begin{figure}[H] 
\centering
\includegraphics[scale = 0.45]{Datenauswertung_P25_Nov2025_IsotropesVariogramm} 
\caption{\small{Schätzung der richtungsunabhängigen 
Semivarianz nach Matheron (links) und Cressie (rechts)
für die $\mathrm{P}_{2.5}$-Konzentrationen von November 2025}}
\label{Abb_Datenauswertung_P25_Nov2025_IsotropesVariogramm}
\end{figure} 

\begin{figure}[H] 
\centering
\includegraphics[scale = 0.45]{Datenauswertung_P25_Nov2025_IsotropesVariogrammAnpassungen} 
\caption{\small{Anpassung (\textcolor{CornflowerBlue}{\rule[1mm]{3mm}{0.5mm}}) 
an verschiedene Kovarianz-Modelle der
geschätzten Semivarianz für die $\mathrm{P}_{2.5}$-Konzentrationen 
von November 2025; eingezeichnet sind zusätzlich
Sill (\textcolor{Green}{\rule[1mm]{3mm}{0.5mm}}),
Range (\textcolor{Brown}{\rule[1mm]{3mm}{0.5mm}}) und
Nugget (\textcolor{LimeGreen}{\rule[1mm]{3mm}{0.5mm}})}}
\label{Abb_Datenauswertung_P25_Nov2025_IsotropesVariogrammAnpassungen}
\end{figure} 
